\documentclass[11pt]{article}
\usepackage[margin=1.1in]{geometry}
\usepackage{times}
\usepackage[T1]{fontenc}
\usepackage{setspace}
\onehalfspacing
\usepackage{parskip}
\pagenumbering{gobble}

\title{There Are No Windows on the Space Station}
\author{Flyxion}
\date{September 2026}

\begin{document}
\maketitle

Every wall on the ring is a display, and every display shows space, rendered at a resolution nobody's eye could distinguish from the article itself, curated slightly, the starfield adjusted for contrast, a nearby nebula given a little more presence than the raw light would earn it on its own. Residents are told this is an upgrade over glass. Glass scratches, glass needs shielding against debris, glass shows you only the fixed slice of sky your window happens to face, while the display can show you any slice, any angle, the whole rotating sphere of it, on request.

The actual windows, the hull's original observation ports, still exist in the schematics, sealed behind the paneling during a refit nobody remembers voting on, justified at the time as a radiation and micrometeorite mitigation measure, technically true, and never reopened afterward because reopening them would have meant admitting the displays were a downgrade wearing an upgrade's marketing copy. A resident who wants the unmediated version, light that traveled untouched from wherever it started and hit a retina without a processing step in between, has to file a request, cite a reason, and wait for a maintenance window, because true windows are now a maintenance category rather than an architectural default.

Nobody experiences this as deprivation, which is the part worth sitting with. The rendered view is better in every measurable way and worse in exactly one, which is that it is a view of space rather than space, a claim so easy to forget once the picture is good enough that most residents, asked, insist they've been looking at the real thing the whole time. Reality did not disappear. It became a premium tier, available on request, filed under maintenance, indistinguishable from the standard package until the day the standard package updates itself and nobody thinks to check whether the light behind it changed too.

\end{document}
